\subsection{No uniform finite list of primes in positive $2$-rank}
\label{sec:finite-prime-avoidance}

\paragraph{Published input (Ozman--Pries).}
Let $\ell\ne p$ be prime, let $g\ge2$, and let $0\le f\le g$, with
$f\ne0$ when $g=2$.  For each irreducible component $Z$ of
$\mathcal M_g^f$, the Raynaud theta divisor of the geometric generic
curve of $Z$ contains no point of order $\ell$.  This is
Ozman--Pries, \emph{Ordinary and almost ordinary Prym varieties},
Asian J. Math. \textbf{23} (2019), 455--478,
Theorem~1.1(2), p.~456, equivalently Theorem~8.3, p.~473.
In particular its hypotheses allow $p=2$, $g=2$, and either $f=1$ or
$f=2$.  The assertion concerns every geometric point of order $\ell$,
not just a selected cyclic cover.

\paragraph{Finding A08 (finite prime avoidance).}
For either $f\in\{1,2\}$ and every finite set $S$ of odd primes, there
exist an integer $m\ge1$ and a smooth projective geometrically connected
curve $C/\mathbf F_{2^m}$ of genus two and $2$-rank $f$ such that
\[
 \Theta_C(\overline{\mathbf F}_2)\cap
 \bigl(J_{C^{(2)}}[\ell](\overline{\mathbf F}_2)\setminus\{0\}\bigr)
 =\varnothing\qquad(\ell\in S).
 \tag{FP.1}
\]
Here $C^{(2)}$ is the relative Frobenius twist, and $\Theta_C$ is the
theta support of the sheaf of locally exact differentials on $C^{(2)}$.
This is the convention of Ozman--Pries, Section~8.1, p.~472.
One can equivalently use the usual Frobenius-semilinear relabelling of
this support; the existence and order of its geometric torsion points
are unchanged.  For $f=2$ the curve is ordinary.

\paragraph{Recovered source scope.}
The ordinary finite-list obstruction already appears in
\texttt{ANABELIAN\_SUCCESSOR}, A.16, whose statement allows arbitrary
characteristic and genus.  Its full recovered proof uses generic
ordinary Pryms and finite cover moduli.  The argument below states the
geometric theta-support conclusion explicitly, gives a scheme-family
spreading proof without coarse-moduli descent, and also treats the
fixed $2$-rank-one stratum.

\begin{proof}
\emph{1. A curve at a geometric generic point.}
The strata $\mathcal M_2^1$ and $\mathcal M_2^2$ are nonempty.  For
example, the smooth projective models of
\[
 y^2+y=x^3+x^{-1},\qquad
 y^2+y=x+x^{-1}+(x-1)^{-1}
\]
have genus two and $2$-ranks one and two, respectively.  For an
Artin--Schreier equation with $r+1$ poles of odd orders $d_j$, the
genus is $\bigl(\sum_j(d_j+1)-2\bigr)/2$ and the $2$-rank is $r$;
see Pries--Zhu, \emph{The $p$-rank stratification of Artin--Schreier
curves}, Ann. Inst. Fourier \textbf{62} (2012), 707--726,
Lemma~2.6, p.~712.

Fix an irreducible component of $\mathcal M_2^f$ over
$\overline{\mathbf F}_2$ and a curve $C_\Omega$ at its geometric generic
point, where $\Omega$ is an algebraically closed field of characteristic
two.  Apply the published input to each $\ell\in S$, on this same
component and this same geometric generic curve.  Thus $C_\Omega$
satisfies all the required avoidance conditions simultaneously.

\emph{2. Spreading the curve while preserving its $2$-rank.}
Choose a point $P\in C_\Omega(\Omega)$ and a projective embedding.
The equations and the point involve finitely many coefficients.
They therefore descend to an integral finitely generated
$\mathbf F_2$-subalgebra $A\subset\Omega$; the precise limit reference
for the scheme and the marked-point morphism is EGA~IV$_3$,
Theorem~8.8.2(ii) and (i), printed p.~28, respectively;
see also the Stacks Project, Lemma~32.10.1(1)--(3),
Tag~\texttt{01ZM}.
After replacing
$T=\operatorname{Spec}A$ by a nonempty affine open, they give a smooth
projective family $\pi:\mathcal C\to T$ of geometrically connected
genus-two curves with a section, whose base change to $\Omega$ is
$(C_\Omega,P)$.  This uses the usual finite-presentation spreading
argument: the projective model is proper, the nonsmooth locus has
closed image missing the generic point, and we remove that image.
The geometric-connectedness locus is constructible by EGA~IV$_3$,
Theorem~9.7.7(ii), p.~79; since it contains the generic point of
the integral base, it contains a nonempty open.  The genus is
preserved by constancy of the Euler characteristic in this smooth
proper family, EGA~III$_2$, Theorem~7.9.4, pp.~76--77.
There is no need to choose a universal curve over a coarse moduli
space or to descend a generic curve to that coarse space's function
field.

We may also ensure that every geometric fibre has $2$-rank $f$.
For completeness, locally trivialize the rank-two bundle
$E=R^1\pi_*\mathcal O_{\mathcal C}$ and let $H$ be the matrix of its
Frobenius-semilinear endomorphism.  On a geometric fibre, the stable
rank of Frobenius is the $2$-rank, and in dimension two the image has
stabilized by the second iterate.  Hence the $2$-rank is
$\operatorname{rank}(H H^{(2)})$, where $H^{(2)}$ denotes entrywise
squaring.  All minors of size $f+1$ vanish at the generic point and
thus vanish identically on the integral base; some minor of size $f$
is nonzero at the generic point.  Removing its zero locus gives the
assertion.  Equivalently one can use the locally closed $p$-rank
strata recalled in Ozman--Pries, Section~2.5, p.~458.

\emph{3. The relative theta incidence is closed.}
Put $\mathcal C^{(2)}=\mathcal C\times_{T,F_T}T$, and let
$F:\mathcal C\to\mathcal C^{(2)}$ be relative Frobenius.
For a smooth relative curve in characteristic two it is finite locally
free of degree two, and
\[
 0\longrightarrow\mathcal O_{\mathcal C^{(2)}}
 \longrightarrow F_*\mathcal O_{\mathcal C}
 \longrightarrow\mathcal B\longrightarrow0
 \tag{FP.2}
\]
has a line-bundle quotient $\mathcal B$, compatible with base change.
These claims can be checked in a relative etale coordinate, where
$F_*\mathcal O$ has basis $1,t$ over the Frobenius-twisted coordinate
ring.  On each fibre, $\mathcal B$ is the sheaf used in the definition
of the Raynaud theta divisor.

Let $\mathcal J=\operatorname{Pic}^0(\mathcal C^{(2)}/T)$, an abelian
scheme.  The chosen section gives a Poincare line bundle $\mathcal P$
on $\mathcal C^{(2)}\times_T\mathcal J$, normalized along that section.
Upper semicontinuity of cohomology for its proper projection to
$\mathcal J$ shows that
\[
 Z_\Theta=
 \{z\in\mathcal J:
 h^0((\mathcal C^{(2)}\times_T\mathcal J)_{\overline z},
       (\operatorname{pr}_1^*\mathcal B\otimes\mathcal P)_{\overline z})>0\}
 \tag{FP.3}
\]
is a closed subset, where $\overline z$ denotes a geometric point
over $z$.  Its geometric fibres over $T$ are
exactly the supports of the relevant theta divisors.  Thus the
argument requires only this closed incidence locus, and does not
require a separate global assertion that a chosen relative theta
divisor is Cartier.  The precise semicontinuity theorem is
Grothendieck, EGA~III$_2$, Theorem~7.7.5(I), printed p.~67;
equivalently, the Stacks Project, Lemma~36.32.1, Tag~\texttt{0BDN}.
Its hypotheses hold because the base is noetherian, the projection
is proper and smooth, and
$\operatorname{pr}_1^*\mathcal B\otimes\mathcal P$ is a line bundle,
hence coherent and flat over $\mathcal J$.

\emph{4. All nonzero torsion points are removed at once.}
For each odd prime $\ell$, the group scheme $\mathcal J[\ell]$ is
finite etale over $T$.  Its zero section is both open and closed, so
\[
 \mathcal J[\ell]^\times:=\mathcal J[\ell]\setminus0_T
\]
is still finite etale over $T$.  The intersection
$Z_\ell=Z_\Theta\cap\mathcal J[\ell]^\times$ is closed in this finite
scheme.  Consequently its image $D_\ell\subset T$ is closed, because
finite morphisms are closed.  A geometric fibre over $t$ contains a
nonzero $\ell$-torsion point on its theta support if and only if
$t\in D_\ell$.

By Step~1, $D_\ell$ misses the generic point of $T$ for every
$\ell\in S$: otherwise its nonempty generic fibre would remain
nonempty after the field extension to $\Omega$.  Therefore
\[
 U=T\setminus\bigcup_{\ell\in S}D_\ell
\]
is a nonempty open subscheme, containing the generic point.
It is of finite type over $\mathbf F_2$, so it has a closed point $u$
with finite residue field $\kappa(u)=\mathbf F_{2^m}$, by the
Nullstellensatz.  The curve $C=\mathcal C_u$ has genus two and
$2$-rank $f$, and satisfies (FP.1) over an algebraic closure of its
residue field.  This proves the claim.
\end{proof}

\paragraph{Ordinary strengthening of A08 (any finite list of odd orders).}
For $f=2$, the finite set $S$ may be replaced by any finite set
$\mathcal N$ of odd integers $n\ge3$: one can require
\[
 \Theta_C(\overline{\mathbf F}_2)\cap
 \bigl(J_{C^{(2)}}[n](\overline{\mathbf F}_2)\setminus\{0\}\bigr)
 =\varnothing\qquad(n\in\mathcal N).
 \tag{FP.4}
\]
In particular there is no uniform finite list of odd prime-power
orders for the sufficient incidence condition in A11.

\begin{proof}
Yang, \emph{Raynaud--Tamagawa theta divisors and new-ordinariness of
ramified coverings of curves}, J. Algebra \textbf{587} (2021),
263--294, Theorem~3.5, applies to arbitrary cyclic degree $m$
prime to the characteristic over a generic pointed stable curve.
For an unpointed smooth curve every ramification divisor $D(j)$
is zero, which is Frobenius stable by Definition~3.3 and
Lemma~3.2.  Hence every connected cyclic etale cover of the
generic genus-two curve, of any odd degree, has ordinary new
factor.  The generic curve itself is ordinary, as also recorded
in Yang's Theorem~1.1.  These statements concern this same
geometric generic curve.

To verify that this excludes odd-order theta points, suppose $L$
has exact odd order $d>1$ on its Raynaud theta support, viewed on
the Frobenius twist.  Choose $e$ with $2e\equiv1\pmod d$ and set
$M=(F^*L)^e$ on the original curve.  Then $M$ has order $d$,
$M^2=F^*L$, and $M^{(2)}=L$; the last equality follows by applying
$F^*$, which is injective on odd-order torsion.
The twisted Frobenius sequence
\[
 0\longrightarrow L\longrightarrow F_*F^*L
 \longrightarrow B\otimes L\longrightarrow0
\]
therefore identifies $H^0(B\otimes L)$ with the kernel of the
semilinear Frobenius arrow
$H^1(C,M)\to H^1(C,M^2)$.
The two spaces are lines, since $C$ has genus two and the line
bundles are nontrivial of degree zero.  This arrow is zero.
In the connected degree-$d$ Kummer cover defined by $M$, it is
one arrow in the nontrivial-character sum of $H^1(\mathcal O)$.
That sum is Frobenius stable, and its Frobenius is not injective,
so its stable rank is smaller than its dimension.  Its new
factor is therefore not ordinary, contradicting Yang's theorem.
Thus the geometric generic theta support contains no nonzero
point of any odd order.

Now use this generic ordinary curve in Steps~2--4 above.  For each
$n\in\mathcal N$, the scheme $\mathcal J[n]\setminus0_T$ is finite
etale.  Its theta incidence has closed image missing the generic
point, because every nonzero element has odd order dividing $n$.
Remove the finitely many images and choose a closed point of the
remaining open.  This gives (FP.4).
\end{proof}

This stronger odd-order assertion is stated only for the ordinary
case $f=2$.  Yang's Section~3.6.2 and Corollary~3.6 retain a prime
degree hypothesis on a generic fixed lower-$p$-rank stratum; no
prime-power extension for $f=1$ is inferred here.

\paragraph{The precise nonuniformity conclusion.}
Taking $S$ to be all odd primes at most an arbitrary real number $B$
shows that no bound depending only on characteristic two, genus two
and the fixed positive $2$-rank can guarantee a nonzero prime-order
theta point below that bound for every curve over finite extensions
of $\mathbf F_2$.  More generally, no fixed finite list of candidate
odd primes works for all those curves.  The extension degree $m$ is
allowed to depend on the list; this proof gives no numerical bound
for $m$ or for the defining equations.

This statement has quantifiers
\[
 \text{for every finite }S\quad\text{there exists }C_S
 \quad\text{avoiding every }\ell\in S.
\]
It does not assert that one finite-field curve avoids every odd prime,
and it does not obstruct a curve-dependent prime of unbounded size.
An infinite intersection of the open avoidance loci need not contain
a closed point.  No analytic hypothesis is used in this deduction.
